\documentclass[12pt,a4paper]{article}

\usepackage[utf8]{inputenc}
\usepackage[vietnamese]{babel}
\usepackage{amsmath,amssymb,amsthm}
\usepackage{geometry}
\usepackage{enumitem}
\usepackage[most]{tcolorbox}
\usepackage{hyperref}
\usepackage{fancyhdr}
\usepackage{tikz}
\usepackage{eso-pic}
\usepackage{xcolor}

\geometry{a4paper, margin=2cm, bottom=2.5cm}
\hypersetup{colorlinks=true, linkcolor=blue!70!black}
\setlist[itemize]{topsep=4pt, itemsep=2pt}
\setlist[enumerate]{topsep=4pt, itemsep=2pt}

\AddToShipoutPictureFG{%
  \AtPageCenter{%
    \begin{tikzpicture}[remember picture, overlay]
      \node[rotate=45, scale=1, opacity=0.06]{\fontsize{60}{72}\selectfont\bfseries\sffamily Đặng Tấn Quí};
    \end{tikzpicture}%
  }%
}

\pagestyle{fancy}
\fancyhf{}
\fancyhead[L]{\small\textit{Toán kinh tế 2}}
\fancyhead[R]{\small\textit{Đặng Tấn Quí}}
\fancyfoot[C]{\thepage}
\fancyfoot[L]{\footnotesize\textcopyright\ Đặng Tấn Quí}
\fancyfoot[R]{\footnotesize SĐT: 0377\,122\,210}
\renewcommand{\headrulewidth}{0.4pt}
\renewcommand{\footrulewidth}{0.4pt}

\fancypagestyle{plain}{\fancyhf{}\fancyfoot[C]{\thepage}\fancyfoot[L]{\footnotesize\textcopyright\ Đặng Tấn Quí}\fancyfoot[R]{\footnotesize SĐT: 0377\,122\,210}\renewcommand{\headrulewidth}{0pt}\renewcommand{\footrulewidth}{0.4pt}}

\newtcolorbox{dinhnghia}[1][]{enhanced, breakable, colback=blue!3!white, colframe=blue!60!black, fonttitle=\bfseries, title={Định nghĩa}, #1}
\newtcolorbox{dinhly}[1][]{enhanced, breakable, colback=green!3!white, colframe=green!50!black, fonttitle=\bfseries, title={Định lý}, #1}
\newtcolorbox{tinhchat}[1][]{enhanced, breakable, colback=yellow!5!white, colframe=orange!50!black, fonttitle=\bfseries, title={Tính chất}, #1}
\newtcolorbox{phuongphap}[1][]{enhanced, breakable, colback=gray!5!white, colframe=gray!50!black, fonttitle=\bfseries, title={Phương pháp}, #1}
\newtcolorbox{luuy}[1][]{enhanced, breakable, colback=red!3!white, colframe=red!50!black, fonttitle=\bfseries, title={Lưu ý}, #1}
\newtcolorbox{congthuc}[1][]{enhanced, breakable, colback=cyan!3!white, colframe=cyan!50!black, fonttitle=\bfseries, title={Công thức}, #1}
\newtcolorbox{vidu}[1][]{enhanced, breakable, colback=violet!3!white, colframe=violet!50!black, fonttitle=\bfseries, title={Ví dụ}, #1}

\begin{document}

\thispagestyle{empty}
\begin{tikzpicture}[remember picture, overlay]
  \fill[blue!80!black] (current page.south west) rectangle (current page.north east);
  \node[anchor=west, white, font=\fontsize{26}{32}\bfseries\selectfont]
    at ([xshift=3cm, yshift=-7.5cm]current page.north west) {TOÁN KINH TẾ 2};
  \node[anchor=west, white, font=\fontsize{18}{24}\selectfont]
    at ([xshift=3cm, yshift=-9cm]current page.north west) {Tích phân, thặng dư, PTVP kinh tế, mô hình Solow};
  \node[anchor=west, white, font=\Large\bfseries] at ([xshift=3cm, yshift=-13cm]current page.north west) {Biên soạn:};
  \node[anchor=west, yellow!80!white, font=\fontsize{22}{28}\bfseries\selectfont] at ([xshift=3cm, yshift=-14.5cm]current page.north west) {ĐẶNG TẤN QUÍ};
  \node[anchor=south east, white, font=\Large\bfseries] at ([xshift=-2cm, yshift=3cm]current page.south east) {\the\year};
\end{tikzpicture}

\newpage
\tableofcontents
\newpage

\section{Tích phân}

\subsection{Nguyên hàm và tích phân xác định}

\begin{dinhnghia}
  $\displaystyle\int f(x)\,dx = F(x) + C$ với $F'(x) = f(x)$. $\displaystyle\int_a^b f(x)\,dx = F(b) - F(a)$.
\end{dinhnghia}

\begin{congthuc}[title={Công thức cơ bản}]
  $\displaystyle\int x^n\,dx = \frac{x^{n+1}}{n+1} + C$ ($n \ne -1$), $\displaystyle\int e^x\,dx = e^x + C$, $\displaystyle\int \frac{1}{x}\,dx = \ln|x| + C$.
\end{congthuc}

\begin{phuongphap}[title={Đổi biến}]
  Đặt $u = g(x)$ $\Rightarrow$ $du = g'(x)dx$. $\displaystyle\int f(g(x))g'(x)\,dx = \int f(u)\,du$.
\end{phuongphap}

\begin{phuongphap}[title={Tích phân từng phần}]
  $\displaystyle\int u\,dv = uv - \int v\,du$.
\end{phuongphap}

\section{Ứng dụng tích phân trong kinh tế}

\subsection{Thặng dư tiêu dùng và thặng dư sản xuất}

\begin{dinhnghia}
  \textbf{Thặng dư tiêu dùng} (CS): phần chênh lệch giữa số tiền người tiêu dùng sẵn sàng trả và số tiền thực tế trả:
  \[
    CS = \int_0^{Q^*} \bigl(P_D(Q) - P^*\bigr)\,dQ = \int_{P^*}^{P_{\max}} Q_D(P)\,dP.
  \]
\end{dinhnghia}

\begin{dinhnghia}
  \textbf{Thặng dư sản xuất} (PS): phần chênh lệch giữa doanh thu thực tế và chi phí tối thiểu sản xuất:
  \[
    PS = \int_0^{Q^*} \bigl(P^* - P_S(Q)\bigr)\,dQ = \int_{P_{\min}}^{P^*} Q_S(P)\,dP.
  \]
\end{dinhnghia}

\begin{vidu}
  Cho hàm cầu $P = 100 - 2Q$, hàm cung $P = 20 + 3Q$. Tìm $P^*$, $Q^*$, CS, PS.
\end{vidu}

\subsection{Dòng tiền và giá trị hiện tại}

\begin{dinhnghia}
  \textbf{Giá trị hiện tại} (PV) của dòng tiền liên tục $f(t)$ trong khoảng $[0,T]$ với lãi suất $r$:
  \[
    PV = \int_0^T f(t)\,e^{-rt}\,dt.
  \]
\end{dinhnghia}

\begin{dinhnghia}
  \textbf{Giá trị tương lai} (FV): $FV = e^{rT} \cdot PV$.
\end{dinhnghia}

\subsection{Hàm cầu từ hàm doanh thu biên}

\begin{phuongphap}
  Nếu biết $MR = R'(Q)$ thì $R(Q) = \displaystyle\int MR\,dQ$, và $P = \dfrac{R}{Q}$.
\end{phuongphap}

\subsection{Tích phân suy rộng trong kinh tế}

\begin{dinhnghia}
  Dòng tiền vô hạn: $PV = \displaystyle\int_0^{+\infty} f(t)\,e^{-rt}\,dt$. Hội tụ khi $f$ có tốc độ tăng chậm hơn $e^{rt}$.
\end{dinhnghia}

\section{Phương trình vi phân trong kinh tế}

\begin{dinhnghia}[title={Mô hình tăng trưởng Solow}]
  $\dot{k} = sf(k) - (n + \delta)k$, $k = K/L$: vốn bình quân.

  Trạng thái dừng $k^*$: $sf(k^*) = (n + \delta)k^*$. Ổn định: $sf'(k^*) < n + \delta$.
\end{dinhnghia}

\begin{dinhnghia}[title={Mô hình Dornbusch (tỷ giá)}]
  Hệ ODE: $\dot{p} = \pi(y - \bar{y})$, $\dot{e} = r - r^*$. Thể hiện overshooting tỷ giá.
\end{dinhnghia}

\begin{phuongphap}[title={PTVP cấp 1 tuyến tính}]
  $\dot{y} + ay = b(t)$. Nghiệm: $y(t) = e^{-at}\left(y_0 + \int_0^t b(s)e^{as}\,ds\right)$.
\end{phuongphap}

\begin{phuongphap}[title={Hệ PTVP tuyến tính}]
  $\dot{\mathbf{x}} = \mathbf{A}\mathbf{x}$. Nghiệm: $\mathbf{x}(t) = e^{\mathbf{A}t}\mathbf{x}_0$. Ổn định: mọi trị riêng $\text{Re}(\lambda) < 0$.
\end{phuongphap}

\section{Phương trình sai phân}

\begin{dinhnghia}[title={Mô hình cobweb (mạng nhện)}]
  $Q_t^S = f(P_{t-1})$, $Q_t^D = g(P_t)$. Cân bằng: $P_t$ hội tụ / phân kỳ / dao động.
\end{dinhnghia}

\begin{phuongphap}[title={Sai phân bậc 2}]
  $y_{t+2} + ay_{t+1} + by_t = c$. Phương trình đặc trưng: $\lambda^2 + a\lambda + b = 0$.

  Ổn định: $|\lambda_1|, |\lambda_2| < 1$.
\end{phuongphap}

\begin{dinhnghia}[title={Mô hình nhân tử--gia tốc (Samuelson)}]
  $Y_t = C_t + I_t + G_0$, $C_t = cY_{t-1}$, $I_t = v(C_t - C_{t-1})$.

  $\Rightarrow$ $Y_{t} - c(1+v)Y_{t-1} + cvY_{t-2} = G_0$. Dao động phụ thuộc $c, v$.
\end{dinhnghia}

\end{document}
